\documentclass[11pt]{article}
\usepackage{amsmath, amssymb, booktabs, geometry}
\geometry{margin=1in}

\title{Biometrical Modelling of Nuclear Family Data:\\
Unified ACDE Variance Component Estimation}
\author{}
\date{}

\begin{document}
\maketitle
\tableofcontents

\section{Introduction}

Biometrical genetics aims to partition phenotypic variation into
genetic and environmental components. Classical models decompose the
phenotypic variance into:

\begin{itemize}
\item Additive genetic effects ($A$)
\item Dominance genetic effects ($D$)
\item Shared environmental effects ($C$)
\item Unique environmental effects ($E$)
\end{itemize}

Nuclear family data provide an alternative to twin designs.
However, identifiability depends strongly on which relatives are observed.

\section{Why Trio Data Are Insufficient}

Consider nuclear family trios (mother, father, child).

Expected covariances:

\begin{center}
\begin{tabular}{lcc}
\toprule
Pair & Expected covariance \\ \midrule
Parent--parent & $C$ \\
Parent--child  & $\tfrac12 A + C$ \\
\bottomrule
\end{tabular}
\end{center}

Dominance effects require two alleles identical by descent.
Parent--child pairs never share dominance deviations:

\[
\text{Cov}_D(\text{parent, child}) = 0
\]

Therefore trio data contain no covariance information about $D$.
This creates confounding:

\[
V_P = A + C + D + E
\]

but only the sum $C + D$ is estimable.

Hence trios allow estimation of AE, ACE or ADE models,
but not the full ACDE model.

\section{Why Adding Siblings Solves Identifiability}

When at least two siblings are observed, a new covariance appears:

\[
\text{Cov}(\text{sibling}_1,\text{sibling}_2)
= \tfrac12 A + C + \tfrac14 D
\]

This additional $\tfrac14 D$ term allows separation of
shared environment and dominance.

\begin{center}
\begin{tabular}{lc}
\toprule
Design & Identifiable components \\ \midrule
Trio & AE, ACE or ADE \\
Nuclear family with siblings & ACDE \\
\bottomrule
\end{tabular}
\end{center}

\section{ACDE Variance Decomposition}

Total phenotypic variance:

\[
V_P = V_A + V_C + V_D + V_E
\]

Interpretation:

\begin{itemize}
\item $V_A$ : additive genetic variance
\item $V_C$ : shared family environment
\item $V_D$ : dominance genetic variance
\item $V_E$ : individual environment + measurement error
\end{itemize}

\section{Expected Covariance Matrix}

For each family we observe:

\[
\mathbf{y}_i = (P_f, P_m, S_1, S_2)^\top
\]

The theoretical covariance matrix is

\[
\Sigma =
\begin{pmatrix}
V_P & V_C & \tfrac12 V_A + V_C & \tfrac12 V_A + V_C \\
V_C & V_P & \tfrac12 V_A + V_C & \tfrac12 V_A + V_C \\
\tfrac12 V_A + V_C & \tfrac12 V_A + V_C & V_P &
\tfrac12 V_A + V_C + \tfrac14 V_D \\
\tfrac12 V_A + V_C & \tfrac12 V_A + V_C &
\tfrac12 V_A + V_C + \tfrac14 V_D & V_P
\end{pmatrix}
\]

Key covariance contributions:

\begin{align}
\text{Parent--child} &= \tfrac12 V_A + V_C \\
\text{Sibling--sibling} &= \tfrac12 V_A + V_C + \tfrac14 V_D \\
\text{Parent--parent} &= V_C
\end{align}

\section{Likelihood Formulation}

Assume phenotypes are multivariate normal:

\[
\mathbf{y}_i \sim \mathcal{N}(\mathbf{0}, \Sigma)
\]

Families are independent, giving log-likelihood:

\[
\ell(\theta) = \sum_{i=1}^N \log f(\mathbf{y}_i|\Sigma(\theta))
\]

For a single family:

\[
\log f(\mathbf{y}_i) =
-\frac12\left[
\log|\Sigma|
+ \mathbf{y}_i^\top \Sigma^{-1} \mathbf{y}_i
+ k\log(2\pi)
\right]
\]

\section{Parameterisation}

To ensure positive variances:

\begin{align}
V_A &= \exp(\theta_A) \\
V_C &= \exp(\theta_C) \\
V_D &= \exp(\theta_D) \\
V_E &= \exp(\theta_E)
\end{align}

Parameters $\theta$ are estimated by maximising the likelihood.

\section{Maximum Likelihood Estimation}

At each iteration:

\begin{enumerate}
\item Construct $\Sigma(\theta)$
\item Compute log-likelihood
\item Update parameters using quasi-Newton optimisation
\end{enumerate}

\section{Derived Quantities}

After estimation:

\[
h^2 = \frac{V_A}{V_P},\quad
c^2 = \frac{V_C}{V_P},\quad
d^2 = \frac{V_D}{V_P},\quad
e^2 = \frac{V_E}{V_P}
\]

\section{Model Comparison}

Nested models are compared using:

\begin{align}
\text{AIC} &= -2\ell + 2k \\
\text{BIC} &= -2\ell + k\log(N)
\end{align}

Likelihood ratio tests:

\[
\chi^2 = -2(\ell_{\text{reduced}} - \ell_{\text{full}})
\]

\section{Conclusion}

Nuclear family data with siblings provide sufficient information to
estimate the full ACDE model using likelihood-based variance component
estimation. The key advantage over trio data is the presence of sibling
covariance, which separates shared environmental and dominance genetic
effects.

\end{document}
